\section{Introduction}
\label{sec:intro}

Deep-research agents---systems that plan, search the web, and synthesize long reports with
citations---are now evaluated almost entirely on \emph{correctness}: whether their claims are
factual and whether citations support them \citep{du2025deepresearchbench,liveresearchbench2025,
mialon2023gaia,wei2025browsecomp}. Correctness is necessary but not sufficient for the use case
these agents are marketed for: helping researchers, analysts, and inventors find ideas that are
not already obvious. That capacity is \emph{creativity}, and it is essentially unmeasured for
this class of system.

We introduce \textbf{CREST} (CREativity of reSearch agenTs), the first released benchmark that
scores deep-research outputs for creativity using the standard construct from the creativity
literature: an output is creative only if it is simultaneously \emph{New}, \emph{Useful}, and
\emph{Surprising} \citep{runco2012standard,boden2004creative,amabile1982cat}. CREST organizes
questions by \emph{task type}---inventing under constraints, counterfactual reasoning, cross-domain
synthesis, meta-creativity, and open-ended generation.

\paragraph{What this paper does and does not report.} An earlier draft of this section announced a
finding that \textbf{no} data in this artifact supports: that an exploration-based research agent is
more creative on cross-domain and counterfactual tasks, that component ablations tie each win to a
specific mechanism, and that those wins are corroborated by human raters. Taking those in turn:
that agent was \emph{never entered}---the released registry contains five off-the-shelf commercial
systems and all $250$ scored records belong to them (\S\ref{sec:experiments}); \emph{no} component
ablation was run, so nothing ties any result to a mechanism; and \emph{no} human study exists, so
there is no human corroboration of anything here. The claims are withdrawn rather than deferred, and
\S\ref{sec:caesar} now presents the agent as a hypothesis and an experimental design rather than a
result.

\paragraph{What CREST did measure.} Across the five commercial entrants, the result is
methodological and negative. Per-axis judge severity spans $3.32/3.20/3.74$ points on a $1$--$10$
scale against a between-entrant spread of $2.06$ points, so \textbf{the leaderboard this benchmark
produces is unidentifiable} and no ranking claim is supportable for any entrant. The
agent$\times$category interaction is \textbf{null} ($F=0.654$, $p=0.837$), so we find no evidence
for the task-dependent dissociation earlier drafts built the paper around. Task type does move
scores more than entrant identity does ($2.92$ vs $2.06$ points), and \emph{Useful} is the binding
axis---but both are statements about the items and the rubric, not about the frontier.

\paragraph{Why we report it this way.} A benchmark paper that discovers its own instrument cannot
rank is more useful released than withheld, and the failure mode is easy to miss: severity dominance
is invisible unless per-judge severity is estimated \emph{and} compared against the effect being
measured. We make no claim about how common that omission is, having run no survey. We
therefore foreground the negative result and treat the exploration agent strictly as future work,
noting that a benchmark whose author's own system had topped it would have invited precisely the
self-serving reading the leaderboard literature warns about \citep{singh2025leaderboard}.

\paragraph{Contributions.}
\begin{enumerate}
  \item \textbf{Result (methodological, negative).} On this panel, judge-severity variance exceeds
    the entire between-entrant signal, so a conventionally-constructed creativity leaderboard for
    deep-research agents is not identifiable; we report per-judge severity, one-way and two-way
    reliability, and the null interaction test (\S\ref{sec:experiments}).
  \item \textbf{Benchmark.} CREST---the first released creativity benchmark for deep-research
    agents: New/Useful/Surprising scoring by task type (fifty questions, ten per type across six
    domains), with dated known-answers and checkable hard rules held judge-side behind a short open
    prompt, and a public release of everything except the run records (\S\ref{sec:crest}). A
    per-item expert baseline ships with each challenge but is \textbf{withheld} from the judge, so
    \emph{Surprising} rests on the judge's own prior. The judge panel is
    \textbf{not} human-validated and uses \textbf{no} frozen shared retrieval index; both are future
    work.
  \item \textbf{An experimental design for the mechanism question, registered nowhere.} A
    single-dissociation ablation protocol that \emph{would} test whether exploration and refinement
    drive category-specific gains, stated in full together with the reasons this run cannot execute
    it (\S\ref{sec:caesar}). It is \textbf{not} in
    \texttt{analysis/preregistration.yaml}---that file holds only the five per-category directions
    and the analysis plan---and it is registered in no external venue.
\end{enumerate}

We state the scope plainly (\S\ref{sec:limitations}): there is no mechanism evidence here, no human
validation, and no identifiable ranking.
